\section{Attribute-Guided Learning Module}

This section describes the attribute-guided learning module added on top of the H-CAST backbone \cite{park2025visually}. The purpose of this module is to use CUB attribute annotations \cite{wah2011cub} as auxiliary supervision and to inject attribute-oriented information into species-level prediction. The module is designed specifically for the fine-grained species branch because species recognition requires subtle local cues, such as bill shape, wing pattern, breast color, crown color, and tail shape.

The proposed module has four main components: attribute-group queries, cross-attention over species-level graph features, group-specific attribute prediction heads, and species-level attribute fusion. Importantly, the module does not use one query or one classifier for each individual CUB attribute. Instead, it operates at the level of active attribute groups. Each active group corresponds to a semantic attribute type, such as \textit{bill shape}, \textit{wing color}, \textit{breast pattern}, or \textit{tail shape}.

\subsection{Attribute-Group Queries}

Let the set of active attribute groups be denoted as
\begin{equation}
    \mathcal{G}_{active}
    =
    \{g_1, g_2, \ldots, g_G\},
\end{equation}
where \(G\) is the number of active groups after applying the attribute mask. For the FullAttr-312 setting, \(G=28\) but for selected attribute subsets, groups with no retained attributes are removed.

The model maintains one learnable query for each active attribute group. Let the learnable attribute-query matrix be
\begin{equation}
    Q^{a}
    =
    [\mathbf{q}_{g_1}, \mathbf{q}_{g_2}, \ldots, \mathbf{q}_{g_G}]
    \in \mathbb{R}^{G \times D},
\end{equation}
where \(D\) is the hidden dimension of the H-CAST backbone. Each query \(\mathbf{q}_{g}\) is associated with one semantic attribute group \(g\), not with one individual attribute. Therefore, the number of attribute queries is equal to the number of active attribute groups, rather than the total number of selected attributes.

For a mini-batch of size \(B\), the query matrix is expanded to
\begin{equation}
    Q^{a}_{B}
    \in
    \mathbb{R}^{B \times G \times D}.
\end{equation}
Before being used in attention, the queries are normalized with a layer normalization module. This normalization helps align the scale of the attribute queries with the H-CAST feature representations.

\subsection{Cross-Attention over Species-Level Graph Features}

The attribute queries attend to graph-node features from the species-level stage of H-CAST. Let
\begin{equation}
    H^{s}
    \in
    \mathbb{R}^{B \times N_s \times D}
\end{equation}
denote the species-level graph-node features, where \(N_s\) is the number of graph nodes at the species-level stage. Let
\begin{equation}
    \mathbf{c}^{s}
    \in
    \mathbb{R}^{B \times D}
\end{equation}
denote the species-level class token.

For the CUB three-level hierarchy, attribute-group queries attend to graph-node features from the species-level stage. In the implementation, when the model predicts three levels, the species-level branch uses the graph nodes from the second H-CAST block and the corresponding species class token. The family and order predictions are produced from later, coarser H-CAST stages.

The cross-attention operation uses attribute-group queries as queries, and species-level graph-node features as keys and values:
\begin{equation}
    H^{a}
    =
    \mathrm{MHA}
    \left(
        Q^{a}_{B},
        H^{s},
        H^{s}
    \right),
\end{equation}
where \(\mathrm{MHA}(\cdot)\) denotes multi-head attention \cite{vaswani2017attention}. The output is
\begin{equation}
    H^{a}
    =
    [\mathbf{h}^{a}_{g_1}, \mathbf{h}^{a}_{g_2}, \ldots, \mathbf{h}^{a}_{g_G}]
    \in
    \mathbb{R}^{B \times G \times D}.
\end{equation}
Each vector \(\mathbf{h}^{a}_{g}\) is an attribute-oriented representation for group \(g\). Intuitively, this operation allows each attribute group to extract visual evidence from the species-level graph nodes. For example, a query corresponding to \textit{bill shape} may learn to attend to regions related to the bird's bill, while a query corresponding to \textit{wing color} may learn to focus on wing-related segments. However, the model is not explicitly supervised with pixel-level attribute localization, so these attention patterns should be interpreted as learned internal representations rather than guaranteed attribute localizations.

A key-padding mask is used in the attention operation to ignore invalid graph nodes introduced by padding. After cross-attention, the attribute-oriented features are normalized and passed through dropout:
\begin{equation}
    H^{a}
    =
    \mathrm{Dropout}
    \left(
        \mathrm{LN}(H^{a})
    \right).
\end{equation}

\subsection{Group-Specific Attribute Prediction}

The attribute-oriented features are supervised by group-specific attribute prediction heads. For each active group \(g \in \mathcal{G}_{active}\), the model defines a separate linear classifier:
\begin{equation}
    f_g^{a}: \mathbb{R}^{D} \rightarrow \mathbb{R}^{C_g},
\end{equation}
where \(C_g\) is the number of retained attribute states in group \(g\). The classifier for group \(g\) receives the corresponding attribute-oriented feature \(\mathbf{h}^{a}_{g}\) and produces logits:
\begin{equation}
    \mathbf{z}^{a}_{g}
    =
    f_g^{a}
    \left(
        \mathrm{LN}(\mathbf{h}^{a}_{g})
    \right).
\end{equation}

This design means that the model does not predict a single 312-dimensional binary attribute vector. Instead, it predicts one retained attribute category for each active attribute group. For example, the classifier for \textit{bill shape} predicts among the retained bill-shape states, while the classifier for \textit{wing color} predicts among the retained wing-color states.

The group-specific attribute predictions are used only as auxiliary outputs during training. Let \(\tilde{t}_{i,g}\) be the remapped target label for sample \(i\) and group \(g\). Let \(\Omega_g\) denote the set of samples whose remapped target for group \(g\) is valid, i.e., not ignored by the selected mask. The set of valid groups in a mini-batch is
\begin{equation}
    \mathcal{G}_{valid}
    =
    \{g \in \mathcal{G}_{active} \mid |\Omega_g| > 0\}.
\end{equation}
The attribute prediction loss is computed as
\begin{equation}
    \mathcal{L}_{attr}
    =
    \frac{1}{|\mathcal{G}_{valid}|}
    \sum_{g \in \mathcal{G}_{valid}}
    \frac{1}{|\Omega_g|}
    \sum_{i \in \Omega_g}
    \mathcal{L}_{CE}
    \left(
        \mathbf{z}^{a}_{i,g},
        \tilde{t}_{i,g}
    \right)
\end{equation}

The role of this auxiliary attribute loss is to encourage each group query to learn features related to its corresponding semantic trait. These attribute predictions are not used as discrete inputs for species classification. Instead, the continuous attribute-oriented features are used in the fusion stage.

\subsection{Species-Level Attribute Fusion}

After obtaining the attribute-oriented features \(H^{a}\), the model aggregates them into a single attribute summary vector. This is done using attention pooling conditioned on the species-level class token. The species class token is used as the query, while the attribute-oriented features are used as keys and values:
\begin{equation}
 \tilde{\mathbf{h}}^{a}_{sum}
=
\mathrm{MHA}
\left(
    \mathbf{c}^{s}_{1},
    H^{a},
    H^{a}
\right),
\qquad
\mathbf{c}^{s}_{1} \in \mathbb{R}^{B \times 1 \times D},
\end{equation}
where \(\mathbf{c}^{s}_{1}\) is the species class token with an added sequence dimension. The output is then squeezed to obtain
\begin{equation}
    \mathbf{h}^{a}_{sum}
    \in
    \mathbb{R}^{B \times D}.
\end{equation}

This attention-pooling step allows the species-level visual representation to select and combine useful attribute-group features. Instead of treating all attribute groups equally, the model can assign different importance to different attribute-oriented representations depending on the input image.

The resulting attribute summary is concatenated with the species-level class token:
\begin{equation}
    \mathbf{h}^{fuse}_{in}
    =
    [
        \mathbf{c}^{s};
        \mathbf{h}^{a}_{sum}
    ]
    \in
    \mathbb{R}^{B \times 2D}.
\end{equation}

The concatenated vector is then passed through a fusion multilayer perceptron:
\begin{equation}
    \mathbf{h}^{s}_{fused}
    =
    \mathrm{MLP}_{fuse}
    \left(
        \mathbf{h}^{fuse}_{in}
    \right),
\end{equation}
where the fusion module maps the \(2D\)-dimensional concatenated representation back to the original hidden dimension \(D\). In the implementation, this fusion module consists of a linear layer, a GELU activation, and layer normalization.

The final species classifier uses the fused species representation:
\begin{equation}
    \mathbf{z}^{s}
    =
    f_s
    \left(
        \mathbf{h}^{s}_{fused}
    \right).
\end{equation}

The attribute-guided fusion is applied to the species-level branch. Family and order predictions are still produced from the original H-CAST hierarchical stages. In the three-level CUB setting, the family classifier uses the family-level H-CAST class token, and the order classifier uses the coarser order-level H-CAST class token:
\begin{equation}
    \mathbf{z}^{f}
    =
    f_f(\mathbf{c}^{f}),
    \qquad
    \mathbf{z}^{o}
    =
    f_o(\mathbf{c}^{o}).
\end{equation}

Therefore, attributes do not directly modify the family or order classifiers. Their primary effect is on species-level classification, where fine-grained attribute cues are expected to be most useful. Family and order predictions may still be indirectly affected during training because the whole network is optimized jointly, but there is no direct attribute fusion into those branches.

\subsection{Summary}

The attribute-guided learning module extends H-CAST by adding attribute-oriented supervision and species-level fusion. It uses one learnable query for each active attribute group, applies cross-attention over species-level graph-node features, predicts retained attribute categories with group-specific classifiers, and fuses the attention-pooled attribute summary with the species class token. This design allows the model to exploit CUB attributes during training while keeping inference image-based. At test time, no manually provided attribute labels are required; the attribute queries and fusion module operate internally using only the image-derived H-CAST features.